\section{Related Work}
\label{sec:related}

\para{Token erasure and implicit vocabulary.}
\citet{feucht2024token} introduced the concept of an implicit vocabulary in LLMs, demonstrating that at the last token position of multi-token words and named entities, linear probes trained to predict nearby tokens show a sharp drop in accuracy in early layers---a phenomenon they term ``token erasure.''
They define an erasure score $\erasure$ and use a greedy segmentation algorithm to extract implicit vocabulary items, recovering approximately 1{,}800 sequences from Wikipedia for Llama-2-7b.
Their analysis reveals that 44.9\% of recovered items are multi-token words or entities, but does not distinguish compositional from non-compositional sequences among the remaining items.
Our work extends this line of research by removing the dependence on pre-trained probes and by systematically characterizing the compositionality spectrum of the recovered items.

\para{Compositionality in neural language models.}
\citet{liu2022representations} studied whether LM representations are locally compositional by training affine probes to predict parent phrase embeddings from constituent embeddings.
They found that human compositionality judgments align weakly with model-internal compositionality scores (best Spearman $\rho = 0.19$ for RoBERTa), suggesting that models do not straightforwardly encode compositionality.
\citet{garcia2021probing} probed for idiomaticity in vector space models and found that contextual representations from BERT partially capture non-compositional semantics.
\citet{singhvi2024shapley} used Shapley interaction indices to measure non-additive interactions between input tokens, showing that autoregressive models encode interactions correlating with syntactic proximity and that both autoregressive and masked models encode nonlinear interactions in idiomatic phrases.
Unlike these works, which start from known non-compositional expressions and probe whether models distinguish them, we start from the model's internal representations and ask what kinds of sequences it treats as unified units.

\para{Multiword expression processing in transformers.}
A growing body of work examines how transformers handle multiword expressions (MWEs).
\citet{nedumpozhimana2022survey} survey methods for detecting and classifying MWEs in transformer models, identifying challenges in distinguishing idiomatic from literal uses.
\citet{zhou2023pier} propose PIER, a model building on BART that creates semantically meaningful representations for potentially idiomatic expressions.
The CLCL approach uses contrastive learning and curriculum learning for non-compositional expression detection~\cite{clcl2023}.
Our work differs in that we do not train a dedicated MWE detector; instead, we analyze the model's existing representations to discover which multi-token sequences it has already learned to treat as vocabulary items.

\para{Internal mechanisms of LLMs.}
\citet{geva2020transformer} showed that transformer feed-forward layers act as key-value memories, storing patterns in keys and associated distributions in values.
\citet{voita2023neurons} categorized neurons in LLMs, identifying ``n-gram neurons'' that activate on specific token sequences, providing complementary evidence that models learn multi-token patterns at the neuron level.
These mechanistic findings suggest that implicit vocabulary items may be supported by dedicated neurons or memory patterns in feed-forward layers, though we leave this connection to future work.

\para{Tokenization and vocabulary design.}
\citet{yang2024rethinking} proposed the Less-is-Better (LiB) tokenizer that learns a vocabulary of subwords, words, and MWEs, showing that incorporating multi-token units into the tokenizer can improve downstream performance.
This connects directly to our finding that models internally develop vocabulary items that a better tokenizer could capture explicitly.
The implicit vocabulary we characterize provides empirical evidence for which multi-token sequences would benefit from explicit tokenization.
